% !TeX root = ../../cubic-quantization.tex
\section{Fermat型の明示量子化}\label{sec:pg-special}
\begin{lemma}\label{lem:pg-fermat-center}
結合代数において
\begin{equation}
[x,y]=\hbar z^2,\qquad [y,z]=\hbar x^2,\qquad [z,x]=\hbar y^2
\end{equation}
が成り立つとする。$S=x^3+y^3+z^3$、$T=\sym(xyz)$ とおくと $S-2\hbar^2T$ は中心元である。
ここで $\sym(xyz)$ は六個の語の平均である。
\end{lemma}
\begin{proof}
関係式と積の交換子を展開すると
\begin{equation}
[S,x]=\hbar(zy^2z-yz^2y)=\hbar[zy,yz]=\hbar^2[zy,x^2]
\end{equation}
を得る。一方、$6T$ の六個の語を関係式で並べ替えると
$2T=xzy+zyx+\hbar x^3$ となる。従って $2[T,x]=[zy,x^2]$ であり、
$[S-2\hbar^2T,x]=0$ を得る。循環置換により $y,z$ との交換子も零である。
\end{proof}
特に
\begin{equation}\label{eq:pg-fermat}
F=\frac13\left(\sum_{\alpha=1}^m u_\alpha^3+
\sum_{s=1}^k\sum_{i=1}^3x_{si}^3\right)
\end{equation}
に対する量子化では、各ブロックで $[x_{s1},x_{s2}]=\hbar x_{s3}^2$ とその循環置換を課し、
異なるブロックと中心変数は可換とする。このとき
\begin{equation}\label{eq:pg-fermat-center}
\Omega=\frac13\left[\sum_{\alpha=1}^m u_\alpha^3+
\sum_{s=1}^k\left\{\sum_{i=1}^3x_{si}^3-2\hbar^2\sym(x_{s1}x_{s2}x_{s3})\right\}\right]
\end{equation}
は全パラメータで中心元であり、$\Omega|_{\hbar=0}=F$ を満たす。
定理\ref{thm:pg-quantization}の適切な開集合上で中心正則超曲面の族を与える。


\subsection{Poisson量子化の明示性}
\begin{proposition}\label{prop:poisson-quantization-explicit}
古典的入力である正規形\eqref{eq:pg-normal}
$F=\sum_{s=1}^k f_s(x_s,u)$に対し、
量子化を定める二次関係式は明示的に
\begin{equation}
 [x_{s2},x_{s3}]
 =\hbar\,\operatorname{sym}_R(\partial_{x_{s1}}f_s),
 \quad
 [x_{s3},x_{s1}]
 =\hbar\,\operatorname{sym}_R(\partial_{x_{s2}}f_s),
 \quad
 [x_{s1},x_{s2}]
 =\hbar\,\operatorname{sym}_R(\partial_{x_{s3}}f_s)
 \label{eq:ci-quantization-relations}
\end{equation}
と書ける。異なるブロックと中心変数の間の交換子は零である。
三次中心元$\Omega_\hbar$は零のZariski近傍で存在し、
有限自由加群間の中心性方程式
$z\mapsto([z,x_{si}])_{s,i}$の核を解いて選べる。
しかし任意の入力に一様な閉公式が得られたとは主張しない。
\end{proposition}
\begin{proof}
二次関係式は各ブロックのcyclic potential
$W_s=x_{s1}x_{s2}x_{s3}-x_{s2}x_{s1}x_{s3}
-\hbar\,\operatorname{sym}_R(f_s)$のcyclic derivativeである。
中心元の存在と、核の計算を用いた代数化は
定理\ref{thm:pg-universal-determinant}および定理\ref{thm:pg-quantization}による。
\end{proof}
\begin{example}[Fermat型]
$[x,y]=\hbar z^2$、$[y,z]=\hbar x^2$、
$[z,x]=\hbar y^2$なら
\begin{equation}
 \Omega=x^3+y^3+z^3
      -2\hbar^2\operatorname{sym}(xyz)
 \label{eq:ci-fermat-center}
\end{equation}
は実際に中心元となる。直接の交換子計算は
補題\ref{lem:pg-fermat-center}にある。
\end{example}
